% ════════════════════════════════════════════════════════════════
%  CORRECTION — DNB BLANC N°2 — Collège Gaston Chaissac, Pouzauges
%  Session du jeudi 9 avril 2026
%  Auteur : M. Belhaj
%  Compilé avec : pdflatex correction-dnb-blanc-2-2026.tex
% ════════════════════════════════════════════════════════════════

\documentclass[a4paper,11pt]{article}

%% ─── Encodage et langue ──────────────────────────────────────
\usepackage[utf8]{inputenc}
\usepackage[T1]{fontenc}
\usepackage[french]{babel}
\usepackage[autolanguage]{numprint}
\newcommand{\np}[1]{\numprint{#1}}

%% ─── Police professionnelle ──────────────────────────────────
\usepackage{fourier}

%% ─── Mise en page ────────────────────────────────────────────
\usepackage[top=2cm, bottom=2cm, left=2cm, right=2cm]{geometry}
\usepackage{fancyhdr}
\pagestyle{fancy}
\fancyhf{}
\fancyhead[L]{\small\textit{Collège Gaston Chaissac — Pouzauges}}
\fancyhead[R]{\small\textit{Correction DNB Blanc n°2 — Avril 2026}}
\cfoot{\thepage}
\renewcommand{\headrulewidth}{0.4pt}
\setlength{\parindent}{0pt}

%% ─── Mathématiques ───────────────────────────────────────────
\usepackage{amsmath, amssymb, amsthm, mathtools}

%% ─── Listes ──────────────────────────────────────────────────
\usepackage{enumitem}

%% ─── Tableaux ────────────────────────────────────────────────
\usepackage{tabularx, array, multirow}

%% ─── Graphiques ──────────────────────────────────────────────
\usepackage{tikz}
\usetikzlibrary{arrows.meta, calc, positioning, angles, quotes, decorations.pathreplacing}
\usepackage{pgfplots}
\pgfplotsset{compat=1.18}

%% ─── Boîtes ──────────────────────────────────────────────────
\usepackage{tcolorbox}
\tcbuselibrary{skins}

%% ─── Divers ──────────────────────────────────────────────────
\usepackage{xcolor}
\usepackage{graphicx}
\definecolor{vertcorrection}{RGB}{22,115,64}
\definecolor{bleutitre}{RGB}{24,66,128}

%% ─── Commandes de mise en forme ──────────────────────────────
\newcommand{\exercice}[3]{%
  \vspace{0.4cm}%
  \noindent\colorbox{bleutitre!15}{%
    \parbox{\dimexpr\linewidth-2\fboxsep}{%
      \color{bleutitre}\Large\bfseries Exercice #1 — #2 \hfill \normalsize\normalfont(#3 points)%
    }%
  }%
  \par\vspace{0.4cm}%
}

\newcommand{\question}[2]{%
  \vspace{0.25cm}%
  \noindent{\bfseries\color{vertcorrection}Question #1}\ \textit{(#2 pt)}\par\vspace{0.15cm}%
}

% Boîte « Réponse » finale
\newcommand{\reponse}[1]{%
  \vspace{0.15cm}%
  \begin{center}
  \fbox{\parbox{0.9\linewidth}{\centering\textbf{Réponse :} #1}}
  \end{center}
  \vspace{0.1cm}%
}

% ══════════════════════════════════════════════════════════════
\begin{document}

% ──────────────────────────────────────────────────────────────
%  PAGE DE GARDE
% ──────────────────────────────────────────────────────────────
\thispagestyle{fancy}

\begin{center}
  {\LARGE\bfseries\color{bleutitre} CORRECTION}\\[0.4cm]
  {\Large\bfseries DNB Blanc n°2 — Mathématiques}\\[0.3cm]
  {\large Collège Gaston Chaissac — Pouzauges}\\[0.2cm]
  {\large Jeudi 9 avril 2026}\\[0.4cm]
  {\itshape M. Belhaj}
\end{center}

\vspace{0.6cm}

\begin{tcolorbox}[colback=gray!10, colframe=gray!50, arc=3pt, boxrule=0.5pt]
\textbf{Barème global :} 24 pts exercices + 2 pts rédaction = 26 pts bruts.\\
\textbf{Note finale :} $\min(\text{total brut},\,20)$ (plafond 20, pas de conversion proportionnelle).

\smallskip

\textbf{Répartition (ordre du bilan élève) :}
\begin{itemize}[leftmargin=1.2cm, itemsep=1pt, topsep=2pt]
  \item Exercice 1 — Automatismes — 6 pts
  \item Exercice 2 — Circuits d'entraînement (PGCD/PPCM) — 3 pts
  \item Exercice 3 — Programme de calcul — 3 pts
  \item Exercice 4 — Jardin botanique (géométrie) — 5 pts
  \item Exercice 5 — Lunettes de soleil (tableur/stats) — 3 pts
  \item Exercice 6 — Fonctions — 4 pts
  \item Rédaction \& justifications — 2 pts (bonus/transparence)
\end{itemize}
\end{tcolorbox}

% ══════════════════════════════════════════════════════════════
%  EXERCICE 1 — Automatismes
% ══════════════════════════════════════════════════════════════
\exercice{1}{Automatismes}{6}

\question{1}{1} Schéma à main levée — ABCD carré, CDE équilatéral (E intérieur), BCF isocèle en F (F extérieur).

\smallskip

On trace un carré ABCD avec le codage habituel :
\begin{itemize}[leftmargin=1cm, itemsep=1pt]
  \item les 4 côtés avec un même codage (petit trait simple) pour indiquer qu'ils ont la même longueur ;
  \item les 4 angles droits aux sommets A, B, C, D.
\end{itemize}

\begin{itemize}[leftmargin=1cm, itemsep=1pt]
  \item Le triangle \textbf{CDE} est équilatéral, E à l'intérieur du carré : ses trois côtés $[CD]$, $[DE]$ et $[EC]$ ont donc la même longueur, égale au côté du carré.
  \item Le triangle \textbf{BCF} a la même taille que CDE (même côté) : $[BC]$, $[BF]$ et $[FC]$ sont aussi égaux au côté du carré. Comme tous les segments de la figure ont la même longueur, on les code d'un seul trait identique.
\end{itemize}

% Figure indicative (schéma à main levée)
\begin{center}
\begin{tikzpicture}[scale=1, thick]
  % Carré ABCD (côté 3)
  \coordinate (A) at (0,3);
  \coordinate (B) at (3,3);
  \coordinate (C) at (3,0);
  \coordinate (D) at (0,0);
  \draw (A)--(B)--(C)--(D)--cycle;
  % Angles droits aux 4 sommets
  \draw (0.25,3) -- (0.25,2.75) -- (0,2.75);
  \draw (2.75,3) -- (2.75,2.75) -- (3,2.75);
  \draw (2.75,0) -- (2.75,0.25) -- (3,0.25);
  \draw (0.25,0) -- (0.25,0.25) -- (0,0.25);

  % E intérieur : triangle équilatéral CDE (côté = 3 = côté du carré)
  \coordinate (E) at (1.5, {3*sin(60)});
  \draw (C)--(E)--(D);

  % F extérieur : triangle équilatéral BCF (côté = 3 = côté du carré, même dimension que CDE)
  \coordinate (F) at ({3 + 3*sin(60)},1.5);
  \draw (B)--(F)--(C);

  % Codage unique : un seul trait sur TOUS les segments (tous égaux au côté du carré)
  \foreach \P/\Q in {A/B, B/C, C/D, D/A, C/E, E/D, B/F, F/C}{
    \path (\P)--(\Q) node[midway, sloped]{\tiny$|$};
  }

  % Labels
  \node[above left]  at (A) {A};
  \node[above right] at (B) {B};
  \node[below right] at (C) {C};
  \node[below left]  at (D) {D};
  \node[below] at (E) {E};
  \node[right] at (F) {F};
\end{tikzpicture}
\end{center}

\reponse{schéma avec codages des 4 côtés du carré (égaux), des angles droits, des 3 côtés du triangle équilatéral CDE, et des deux côtés égaux $[FB]$ et $[FC]$ du triangle isocèle BCF.}

\question{2}{1} Périmètre d'un cercle de rayon 3~cm.

Le périmètre d'un cercle de rayon $r$ est $P = 2\pi r$.
\[
P = 2 \times \pi \times 3 = 6\pi \approx 18{,}85 \text{ cm}
\]

\reponse{$\boxed{P = 6\pi \text{ cm} \approx 18{,}85 \text{ cm}}$}

\question{3}{1} Calculer $77{,}6 \div 10$.

Diviser par 10 revient à décaler la virgule d'un rang vers la gauche.

\reponse{$\boxed{77{,}6 \div 10 = 7{,}76}$}

\question{4}{1} Total des ventes sur la semaine.

On lit sur le diagramme : Lundi 33, Mardi 49, Mercredi 48, Jeudi 30, Vendredi 32.
\[
33 + 49 + 48 + 30 + 32 = 192
\]

\reponse{$\boxed{\text{Total des ventes} = 192}$}

\question{5}{1} Densité de population (\np{49298} habitants, 302~km$^2$).

La densité est le nombre d'habitants par km$^2$ :
\[
d = \dfrac{\np{49298}}{302} \approx 163{,}24 \text{ hab/km}^2
\]

\reponse{$\boxed{d \approx 163 \text{ hab/km}^2}$}

\question{6}{1} Script Scratch — pentagone régulier de côté 147 pixels.

Un pentagone régulier : on tourne de $\dfrac{360^\circ}{5} = 72^\circ$ à chaque sommet.

\begin{center}
\fbox{\parbox{0.7\linewidth}{\ttfamily
Quand (drapeau vert) est cliqué\\
Répéter 5 fois\\
\hspace*{1cm}Avancer de 147\\
\hspace*{1cm}Tourner de 72 degrés\\
Fin
}}
\end{center}

\reponse{boucle « répéter 5 fois » avec « avancer de 147 » puis « tourner de 72 degrés ».}

% ══════════════════════════════════════════════════════════════
%  EXERCICE 2 — Circuits d'entraînement
% ══════════════════════════════════════════════════════════════
\exercice{2}{Circuits d'entraînement (PGCD/PPCM)}{3}

\question{1}{0,5} Montrer que le circuit 1 dure 280~s et le circuit 2 dure 350~s.

\textbf{Circuit 1 :} 5 exercices de 40~s + 5 repos de 16~s (chaque exercice est suivi d'un repos, y compris après le 5\ieme{} pour revenir au départ).
\[
5 \times 40 + 5 \times 16 = 200 + 80 = 280 \text{ s}
\]

\textbf{Circuit 2 :} 10 exercices de 30~s + 10 repos de 5~s.
\[
10 \times 30 + 10 \times 5 = 300 + 50 = 350 \text{ s}
\]

\reponse{$\boxed{\text{Circuit 1} = 280 \text{ s} \;;\; \text{Circuit 2} = 350 \text{ s}}$}

\question{2}{0,5} Décomposition en produit de facteurs premiers de 280 et 350.

\begin{align*}
280 &= 2 \times 140 = 2 \times 2 \times 70 = 2 \times 2 \times 2 \times 35 = 2 \times 2 \times 2 \times 5 \times 7 \\
350 &= 2 \times 175 = 2 \times 5 \times 35 = 2 \times 5 \times 5 \times 7
\end{align*}

\reponse{$\boxed{280 = 2^3 \times 5 \times 7 \qquad 350 = 2 \times 5^2 \times 7}$}

\question{3.a}{1} Position à \np{2800}~s.

\textbf{Camille :} un tour dure 280~s. Or $\np{2800} = 280 \times 10$, donc à \np{2800}~s Camille termine exactement son 10\ieme{} tour du circuit 1 : elle est de retour au départ.

\textbf{Dominique :} un tour du circuit 2 dure 350~s.
\[
\np{2800} \div 350 = 8
\]
Donc à \np{2800}~s, Dominique termine exactement son 8\ieme{} tour du circuit 2. Comme \np{2800} est un multiple de 350, Dominique est également \textbf{au départ du circuit 2}.

\reponse{$\boxed{\text{Camille et Dominique sont tous les deux au départ de leur circuit à } t = \np{2800} \text{ s}}$}

\question{3.b}{1} Première rencontre au départ (PPCM).

On cherche le \textbf{PPCM} de 280 et 350. À partir des décompositions :
\begin{itemize}[leftmargin=1cm, itemsep=1pt]
  \item $280 = 2^3 \times 5 \times 7$
  \item $350 = 2 \times 5^2 \times 7$
\end{itemize}
\[
\text{PPCM}(280\,;\,350) = 2^3 \times 5^2 \times 7 = 8 \times 25 \times 7 = \np{1400} \text{ s}
\]

\textbf{Conversion :} $\np{1400} \div 60 = 23$ minutes, reste $\np{1400} - 23 \times 60 = \np{1400} - \np{1380} = 20$~s.

\reponse{$\boxed{\np{1400} \text{ s} = 23 \text{ min } 20 \text{ s}}$}

% ══════════════════════════════════════════════════════════════
%  EXERCICE 3 — Programme de calcul
% ══════════════════════════════════════════════════════════════
\exercice{3}{Programme de calcul}{3}

\textit{Programme : choisir un nombre $\to$ le mettre au carré $\to$ soustraire le triple du nombre de départ $\to$ soustraire 4.}

\question{1}{0,5} Test avec $x = 5$.
\[
5^2 - 3 \times 5 - 4 = 25 - 15 - 4 = 6
\]

\reponse{$\boxed{\text{Le résultat est bien } 6}$}

\question{2}{0,5} Expression en fonction de $x$.

Avec un nombre de départ $x$, le programme donne :
\[
x^2 - 3x - 4
\]

\reponse{$\boxed{x^2 - 3x - 4}$}

\question{3}{0,5} Vérifier que $x^2 - 3x - 4 = (x+1)(x-4)$.

On développe $(x+1)(x-4)$ :
\begin{align*}
(x+1)(x-4) &= x \times x + x \times (-4) + 1 \times x + 1 \times (-4) \\
           &= x^2 - 4x + x - 4 = x^2 - 3x - 4
\end{align*}

\reponse{$\boxed{(x+1)(x-4) = x^2 - 3x - 4}$ \ \checkmark}

\question{4}{0,5} Pour quels $x$ le programme donne-t-il 0 ?

On résout $(x+1)(x-4) = 0$. Un produit de facteurs est nul si et seulement si l'un des facteurs est nul :
\begin{itemize}[leftmargin=1cm, itemsep=1pt]
  \item $x + 1 = 0 \iff x = -1$
  \item $x - 4 = 0 \iff x = 4$
\end{itemize}

\reponse{$\boxed{x = -1 \text{ ou } x = 4}$}

\question{5}{1} Script Scratch — compléter les lignes 4 et 6.

Le programme calcule $x^2 - 3x - 4$.

\begin{itemize}[leftmargin=1cm, itemsep=2pt]
  \item \textbf{Ligne 4} : \texttt{mettre y à (x * x)} \quad (le carré de $x$)
  \item \textbf{Ligne 6} : \texttt{mettre Résultat à (y - z - 4)} \quad (avec $z = 3x$)
\end{itemize}

\reponse{Ligne 4 : \texttt{mettre y à (x * x)} \quad ;\quad Ligne 6 : \texttt{mettre Résultat à (y - z - 4)}}

% ══════════════════════════════════════════════════════════════
%  EXERCICE 4 — Jardin botanique
% ══════════════════════════════════════════════════════════════
\exercice{4}{Jardin botanique (géométrie)}{5}

\textit{Données :} ABCD quadrilatère ; $AB = 500$~m, $BE = 250$~m, $DE = 750$~m ; $[AC]$ et $[BD]$ se coupent en E ; angles droits en A ($\widehat{DAB}$), en D ($\widehat{ADC}$) et en E (diagonales perpendiculaires).

% Reproduction du schéma du sujet original
\begin{center}
\begin{tikzpicture}[scale=0.85, thick]
  \coordinate (A) at (0,4.33);
  \coordinate (B) at (2.5,4.33);
  \coordinate (C) at (7.5,0);
  \coordinate (D) at (0,0);
  \coordinate (E) at (1.875,3.248);

  % Remplissage gris
  \fill[gray!25] (A) -- (B) -- (E) -- cycle;
  \fill[gray!25] (D) -- (E) -- (C) -- cycle;

  % Quadrilatère
  \draw (A) -- (B) -- (C) -- (D) -- cycle;

  % Diagonales (pointillés)
  \draw[densely dotted, line width=1.2pt] (A) -- (C);
  \draw[densely dotted, line width=1.2pt] (B) -- (D);

  % Angles droits
  \draw (0.25,4.33) -- (0.25,4.08) -- (0,4.08);
  \draw (0.25,0) -- (0.25,0.25) -- (0,0.25);
  \begin{scope}[shift={(E)}, rotate=-71.6]
    \draw (0.2,0) -- (0.2,0.2) -- (0,0.2);
  \end{scope}

  % Labels
  \node[above left]  at (A) {A};
  \node[above right] at (B) {B};
  \node[below right] at (C) {C};
  \node[below left]  at (D) {D};
  \node[left=4pt]    at (E) {E};
\end{tikzpicture}
\end{center}

\question{1}{0,5} Longueur du segment $[DB]$.

E appartient à $[BD]$, donc $DB = BE + ED$.
\[
DB = 250 + 750 = \np{1000} \text{ m}
\]

\reponse{$\boxed{DB = \np{1000} \text{ m}}$}

\question{2}{1} Longueur $[AD]$ par Pythagore dans ABD rectangle en A.

Le triangle ABD est rectangle en A. D'après le \textbf{théorème de Pythagore} :
\begin{align*}
BD^2 &= AB^2 + AD^2 \\
\np{1000}^2 &= 500^2 + AD^2 \\
\np{1000000} &= \np{250000} + AD^2 \\
AD^2 &= \np{750000} \\
AD &= \sqrt{\np{750000}} \approx 866{,}025 \text{ m}
\end{align*}

\reponse{$\boxed{AD \approx 866 \text{ m}}$}

\question{3.a}{0,5} Sinus de l'angle $\widehat{EAB}$.

Dans le triangle ABE, l'angle en E est droit (les diagonales se coupent perpendiculairement en E). Le triangle ABE est donc rectangle en E.

Pour l'angle $\widehat{EAB}$ :
\begin{itemize}[leftmargin=1cm, itemsep=1pt]
  \item côté opposé : $[BE]$ (face à l'angle en A) ;
  \item hypoténuse : $[AB]$ (face à l'angle droit en E).
\end{itemize}

\[
\sin(\widehat{EAB}) = \dfrac{BE}{AB} = \dfrac{250}{500} = \dfrac{1}{2} = 0{,}5
\]

\reponse{$\boxed{\sin(\widehat{EAB}) = 0{,}5}$}

\question{3.b}{0,5} Mesure de l'angle $\widehat{EAB}$.

On cherche $\widehat{EAB}$ tel que $\sin(\widehat{EAB}) = 0{,}5$. À la calculatrice :
\[
\widehat{EAB} = \arcsin(0{,}5) = 30^\circ
\]

\reponse{$\boxed{\widehat{EAB} = 30^\circ}$}

\question{4.a}{0,5} Montrer que $(AB) \parallel (DC)$.

On sait que :
\begin{itemize}[leftmargin=1cm, itemsep=1pt]
  \item $(AB) \perp (AD)$ (car ABD est rectangle en A) ;
  \item $(DC) \perp (AD)$ (car l'angle $\widehat{ADC}$ est droit dans le quadrilatère).
\end{itemize}

Deux droites perpendiculaires à une même troisième droite sont parallèles entre elles.

\reponse{$\boxed{(AB) \parallel (DC)}$}

\question{4.b}{1} Montrer que $CD = \np{1500}$~m (Thalès).

Les droites $(AC)$ et $(BD)$ se coupent en E. D'après la question 4.a, les droites $(AB)$ et $(CD)$ sont parallèles. On applique le \textbf{théorème de Thalès} dans cette configuration (triangles EAB et ECD, en « papillon » avec E sommet commun) :
\[
\dfrac{EA}{EC} = \dfrac{EB}{ED} = \dfrac{AB}{CD}
\]

On utilise le rapport connu :
\[
\dfrac{EB}{ED} = \dfrac{250}{750} = \dfrac{1}{3}
\]

Donc :
\[
\dfrac{AB}{CD} = \dfrac{1}{3} \iff CD = 3 \times AB = 3 \times 500 = \np{1500} \text{ m}
\]

\reponse{$\boxed{CD = \np{1500} \text{ m}}$}

\question{5}{1} Temps du tour du jardin à 1,1~m/s ; est-il inférieur à 1~h ?

On calcule le périmètre du quadrilatère ABCD :
\[
P = AB + BC + CD + DA = 500 + \np{1323} + \np{1500} + 866 = \np{4189} \text{ m}
\]

Le piéton marche à 1,1~m/s. Le temps $t$ est :
\[
t = \dfrac{\np{4189}}{1{,}1} \approx \np{3808}{,}2 \text{ s}
\]

Or 1~heure $= \np{3600}$~s. Donc $t \approx \np{3808}$~s $>$ \np{3600}~s.

\reponse{$\boxed{\text{Non, le temps ($\approx$ \np{3808}~s $\approx$ 1~h 3~min) est supérieur à 1~heure.}}$}

% ══════════════════════════════════════════════════════════════
%  EXERCICE 5 — Lunettes de soleil (ordre app : Ex5 = Lunettes)
% ══════════════════════════════════════════════════════════════
\exercice{5}{Lunettes de soleil (tableur / stats)}{3}

\textit{Données :} prix à l'unité des 5 modèles : 75~€ ; 100~€ ; 110~€ ; 140~€ ; 160~€. Paires vendues : \np{1200} ; 950 ; 875 ; 250 ; 300.

\question{1}{0,5} Étendue des prix.

L'étendue = plus grande valeur $-$ plus petite valeur.
\begin{itemize}[leftmargin=1cm, itemsep=1pt]
  \item Prix maximum : 160~€
  \item Prix minimum : 75~€
\end{itemize}
\[
\text{Étendue} = 160 - 75 = 85 \text{ €}
\]

\reponse{$\boxed{\text{Étendue} = 85 \text{ €}}$}

\question{2.a}{0,5} Formule en G2 (somme des paires vendues, plage B2:F2).

\reponse{$\boxed{\texttt{=SOMME(B2:F2)}}$}

\question{2.b}{0,5} Nombre total de paires vendues.
\[
\np{1200} + 950 + 875 + 250 + 300 = \np{3575} \text{ paires}
\]

\reponse{$\boxed{\text{Total} = \np{3575} \text{ paires}}$}

\question{3.a}{1} Montant total des ventes.

Pour chaque modèle, prix $\times$ quantité, puis on additionne :
\begin{itemize}[leftmargin=1cm, itemsep=1pt]
  \item Modèle 1 : $\np{1200} \times 75 = \np{90000}$ €
  \item Modèle 2 : $950 \times 100 = \np{95000}$ €
  \item Modèle 3 : $875 \times 110 = \np{96250}$ €
  \item Modèle 4 : $250 \times 140 = \np{35000}$ €
  \item Modèle 5 : $300 \times 160 = \np{48000}$ €
\end{itemize}
\[
\np{90000} + \np{95000} + \np{96250} + \np{35000} + \np{48000} = \np{364250} \text{ €}
\]

\reponse{$\boxed{\text{Montant total} = \np{364250} \text{ €}}$}

\question{3.b}{0,5} Prix moyen d'une paire de lunettes (arrondi au centime).

Le prix moyen est le montant total divisé par le nombre total de paires vendues :
\[
\bar{p} = \dfrac{\np{364250}}{\np{3575}} \approx 101{,}8881\ldots \text{ €}
\]

Arrondi au centime près :

\reponse{$\boxed{\bar{p} \approx 101{,}89 \text{ €}}$}

% ══════════════════════════════════════════════════════════════
%  EXERCICE 6 — Fonctions (ordre app : Ex6 = Fonctions)
% ══════════════════════════════════════════════════════════════
\exercice{6}{Fonctions}{4}

\textit{Programme A :} choisir un nombre $\to$ d'un côté « Ajouter 3 » ; de l'autre « Soustraire 4 » $\to$ multiplier les deux résultats. On admet $f(x) = (x+3)(x-4)$.

\question{1.a}{0,5} Test avec $x = -8$.
\begin{itemize}[leftmargin=1cm, itemsep=1pt]
  \item $-8 + 3 = -5$
  \item $-8 - 4 = -12$
  \item Produit : $(-5) \times (-12) = 60$
\end{itemize}

\reponse{$\boxed{\text{Le résultat est bien } 60}$}

\question{1.b}{0,5} Résoudre $(x+3)(x-4) = 0$.

Un produit est nul si et seulement si l'un de ses facteurs est nul :
\begin{itemize}[leftmargin=1cm, itemsep=1pt]
  \item $x + 3 = 0 \iff x = -3$
  \item $x - 4 = 0 \iff x = 4$
\end{itemize}

\reponse{$\boxed{x = -3 \text{ ou } x = 4 \text{ donnent un résultat nul.}}$}

\question{2.a}{0,5} Forme développée de $f(x)$.
\[
(x+3)(x-4) = x^2 - 4x + 3x - 12 = x^2 - x - 12
\]

\reponse{$\boxed{f(x) = x^2 - x - 12}$}

\question{2.b}{0,5} Calcul de $f\!\left(\dfrac{1}{2}\right)$.
\[
f\!\left(\dfrac{1}{2}\right) = \left(\dfrac{1}{2}\right)^2 - \dfrac{1}{2} - 12 = \dfrac{1}{4} - \dfrac{1}{2} - 12
\]

Avec un dénominateur commun 4 :
\[
f\!\left(\dfrac{1}{2}\right) = \dfrac{1}{4} - \dfrac{2}{4} - \dfrac{48}{4} = \dfrac{1 - 2 - 48}{4} = -\dfrac{49}{4} = -12{,}25
\]

\reponse{$\boxed{f\!\left(\tfrac{1}{2}\right) = -\dfrac{49}{4} = -12{,}25}$}

\question{2.c}{0,5} Antécédents de $-6$ par $f$ (graphiquement).

On cherche les $x$ tels que $f(x) = -6$. On trace la droite horizontale d'équation $y = -6$ et on lit les abscisses des points d'intersection avec $\mathcal{C}_f$.

\textit{Vérification algébrique (non demandée) :}
\[
x^2 - x - 12 = -6 \iff x^2 - x - 6 = 0 \iff (x+2)(x-3) = 0
\]
d'où $x = -2$ ou $x = 3$. Graphiquement on lit bien deux antécédents : environ $-2$ et $3$.

\reponse{$\boxed{\text{Les antécédents de } -6 \text{ sont } x = -2 \text{ et } x = 3.}$}

\question{3.a}{0,5} Formule tableur en B2 pour $g(x) = 3x - 7$.

En B2, on calcule $g$ appliqué à la valeur de A2.

\reponse{$\boxed{\texttt{=3*A2-7}}$}

\question{3.b}{0,5} Tracé de la représentation graphique de $g$.

$g$ est une fonction affine : sa représentation est une droite. Deux points du tableau suffisent, par exemple :
\begin{itemize}[leftmargin=1cm, itemsep=1pt]
  \item $(0\,;\,-7)$ (ordonnée à l'origine)
  \item $(4\,;\,5)$
\end{itemize}

On trace la droite passant par ces deux points (coefficient directeur 3).

\reponse{droite passant par $(0\,;\,-7)$ et $(4\,;\,5)$ (coefficient directeur 3).}

\question{3.c}{0,5} Nombres ayant la même image par $f$ et $g$ (graphique).

On cherche les $x$ tels que $f(x) = g(x)$. Graphiquement, ce sont les abscisses des points d'intersection de $\mathcal{C}_f$ (parabole bleue) et de la droite tracée pour $g$.

\textit{Vérification algébrique (non demandée) :}
\[
x^2 - x - 12 = 3x - 7 \iff x^2 - 4x - 5 = 0 \iff (x+1)(x-5) = 0
\]
d'où $x = -1$ ou $x = 5$. On lit graphiquement les deux points d'intersection d'abscisses $-1$ et $5$.

\reponse{$\boxed{\text{Les nombres qui ont la même image par } f \text{ et } g \text{ sont } x = -1 \text{ et } x = 5.}$}

% ══════════════════════════════════════════════════════════════
%  RÉCAPITULATIF BARÈME
% ══════════════════════════════════════════════════════════════
\vspace{0.6cm}

\begin{center}
{\large\bfseries\color{bleutitre} Récapitulatif du barème}
\end{center}

\begin{center}
\renewcommand{\arraystretch}{1.3}
\begin{tabular}{|l|l|c|}
\hline
\textbf{Exercice} & \textbf{Thème} & \textbf{Points} \\ \hline
1 & Automatismes & 6 \\ \hline
2 & Circuits (PGCD/PPCM) & 3 \\ \hline
3 & Programme de calcul & 3 \\ \hline
4 & Jardin botanique & 5 \\ \hline
5 & Lunettes de soleil (tableur/stats) & 3 \\ \hline
6 & Fonctions & 4 \\ \hline
\multicolumn{2}{|r|}{\textbf{Sous-total exercices}} & \textbf{24} \\ \hline
\multicolumn{2}{|l|}{Rédaction \& justifications (bonus/transparence)} & 2 \\ \hline
\multicolumn{2}{|r|}{\textbf{Total brut}} & \textbf{26} \\ \hline
\multicolumn{2}{|l|}{\textbf{Note finale} — $\min(\text{total brut},\,20)$, plafonnée à 20} & \textbf{/20} \\ \hline
\end{tabular}
\end{center}

\vspace{0.4cm}

\begin{center}
\textit{Fin de la correction — DNB Blanc n°2 — Collège Gaston Chaissac, avril 2026 — M. Belhaj.}
\end{center}

\end{document}
